\documentclass[12pt,a4paper,oneside]{article}
\usepackage[brazil]{babel}
\usepackage[utf8]{inputenc}
\usepackage{olymp}

\setcounter{problem}{1}

\begin{document}

\begin{problem}{Ângulo interno}{angulo}{2}

Num polígono regular, todos os ângulos internos tem o mesmo valor. Por exemplo, num triângulo isósceles todos os ângulos valem 60$^{\circ}$. Num quadrado todos são 90$^{\circ}$. E num pentágono regular todos são 108$^{\circ}$.

É possível calcular os ângulos internos de um polígono regular de $N$ lados pela fórmula abaixo:

$$\text{ângulo} = \frac{180(N-2)}{N}$$

Faça um programa que calcula o valor dos ângulos internos de vários polígonos.


%\Notes
%
%Observações, se tiver.

\InputFile

A entrada contém dois valores inteiros, $A$ e $B$, que indicam o menor e o maior número de lados de polígono que seu programa deve tratar. Restrição: $3 \leq A < B \leq 100$.

\OutputFile

Seu programa deve escrever o valor do ângulo interno dos polígonos regulares de $A$ lados, $A+1$ lados, $\dots$, $B$ lados. Cada valor deve ser escrito em uma linha, formatado com 1 casa decimal.

% \Notes
% xxx

\Example

\begin{example}
\exmp{
3 5
}{
60.0
90.0
108.0
}%
\end{example}

\begin{example}
\exmp{
5 10
}{
108.0
120.0
128.6
135.0
140.0
144.0
}%
\end{example}


%Comentários sobre o exemplo, se necessário.
\small
No segundo exemplo são mostrados os ângulos internos dos polígonos de 5, 6, 7, 8, 9 e 10 lados

\end{problem}

\end{document}
